
Code generation models optimized solely for execution correctness produce code that is 19\% slower than human-written alternatives despite passing all tests. This efficiency gap motivates multi-objective optimization approaches that incorporate auxiliary metrics beyond functional correctness. However, existing methods adopt proxy metrics---structural similarity (CodeBLEU), runtime efficiency, style conformity---without pre-validating measurement reliability. This work presents a four-stage validation pipeline to test candidate proxies before reinforcement learning optimization. Stage 1 (measurement reliability, validated via proof-of-concept) evaluates coefficient of variation (CV $\leq 5$\%), effect size (Cohen's d $\geq 0.8)$, and cross-platform stability (Spearman $\rho \geq 0.8)$. Stages 2-4 (conditional independence, cross-domain generalization, optimization constraints) are designed for future validation. Proof-of-concept evaluation using synthetic measurements demonstrates that structural similarity (CodeBLEU) achieves exceptional reliability (CV=1.39\%, Cohen's d=4.51, Spearman $\rho$=0.949), while runtime efficiency measurements using a wall-clock noise model marginally exceed the CV threshold (6.22\% vs 5.0\%). Literature evidence suggests hardware instruction counting via CPU performance counters would achieve lower variance (CV ~2-3\%), though empirical confirmation is needed. The framework employs a scoped gate design where partial validation ($\geq 1$ proxy passing all criteria) constitutes progression, preventing all-or-nothing failure. This work establishes construct validation as a prerequisite for proxy-based optimization, demonstrating that different quality dimensions have distinct measurement reliability profiles requiring independent testing.
